% !TEX program = lualatex
\ifdefined\pg\else\IfFileExists{tex/avulso.tex}{% Compilação AVULSA de um fragmento (volN/.../<stem>.tex) — p. ex. pelo botão do editor,
% de dentro da pasta do nó.  O fragmento começa com
%
%   \ifdefined\pg\else\IfFileExists{tex/avulso.tex}{% Compilação AVULSA de um fragmento (volN/.../<stem>.tex) — p. ex. pelo botão do editor,
% de dentro da pasta do nó.  O fragmento começa com
%
%   \ifdefined\pg\else\IfFileExists{tex/avulso.tex}{% Compilação AVULSA de um fragmento (volN/.../<stem>.tex) — p. ex. pelo botão do editor,
% de dentro da pasta do nó.  O fragmento começa com
%
%   \ifdefined\pg\else\IfFileExists{tex/avulso.tex}{\input{tex/avulso}}{\input{../../tex/avulso}}\fi
%
% e termina com  \csname fimavulso\endcsname .  Dentro do livro ou do _preview, \pg já existe
% e nada disto roda; \fimavulso fica indefinido e o \csname vira \relax.
%
% Motor: o projeto é lualatex (fontspec/polyglossia).  Se o editor chamar pdflatex, cai
% num preâmbulo equivalente com inputenc/babel — serve para conferir, não para o PDF final.
\documentclass[11pt,a5paper]{book}
\ifdefined\directlua
  \usepackage{fontspec}
  \usepackage{polyglossia}
  \setmainlanguage{portuguese}
\else\ifdefined\XeTeXrevision
  \usepackage{fontspec}
  \usepackage{polyglossia}
  \setmainlanguage{portuguese}
\else
  \usepackage[utf8]{inputenc}
  \usepackage[T1]{fontenc}
  \usepackage{lmodern}
  \usepackage[portuguese]{babel}
  \DeclareUnicodeCharacter{2192}{\ensuremath{\rightarrow}}
\fi\fi
% raiz do repo: compilado da raiz ou de volN/<parte>/ (dois níveis abaixo)
\IfFileExists{tex/preambulo.tex}{\input{tex/preambulo}}{\input{../../tex/preambulo}}
\def\fimavulso{\end{document}}
\begin{document}
}{% Compilação AVULSA de um fragmento (volN/.../<stem>.tex) — p. ex. pelo botão do editor,
% de dentro da pasta do nó.  O fragmento começa com
%
%   \ifdefined\pg\else\IfFileExists{tex/avulso.tex}{\input{tex/avulso}}{\input{../../tex/avulso}}\fi
%
% e termina com  \csname fimavulso\endcsname .  Dentro do livro ou do _preview, \pg já existe
% e nada disto roda; \fimavulso fica indefinido e o \csname vira \relax.
%
% Motor: o projeto é lualatex (fontspec/polyglossia).  Se o editor chamar pdflatex, cai
% num preâmbulo equivalente com inputenc/babel — serve para conferir, não para o PDF final.
\documentclass[11pt,a5paper]{book}
\ifdefined\directlua
  \usepackage{fontspec}
  \usepackage{polyglossia}
  \setmainlanguage{portuguese}
\else\ifdefined\XeTeXrevision
  \usepackage{fontspec}
  \usepackage{polyglossia}
  \setmainlanguage{portuguese}
\else
  \usepackage[utf8]{inputenc}
  \usepackage[T1]{fontenc}
  \usepackage{lmodern}
  \usepackage[portuguese]{babel}
  \DeclareUnicodeCharacter{2192}{\ensuremath{\rightarrow}}
\fi\fi
% raiz do repo: compilado da raiz ou de volN/<parte>/ (dois níveis abaixo)
\IfFileExists{tex/preambulo.tex}{\input{tex/preambulo}}{\input{../../tex/preambulo}}
\def\fimavulso{\end{document}}
\begin{document}
}\fi
%
% e termina com  \csname fimavulso\endcsname .  Dentro do livro ou do _preview, \pg já existe
% e nada disto roda; \fimavulso fica indefinido e o \csname vira \relax.
%
% Motor: o projeto é lualatex (fontspec/polyglossia).  Se o editor chamar pdflatex, cai
% num preâmbulo equivalente com inputenc/babel — serve para conferir, não para o PDF final.
\documentclass[11pt,a5paper]{book}
\ifdefined\directlua
  \usepackage{fontspec}
  \usepackage{polyglossia}
  \setmainlanguage{portuguese}
\else\ifdefined\XeTeXrevision
  \usepackage{fontspec}
  \usepackage{polyglossia}
  \setmainlanguage{portuguese}
\else
  \usepackage[utf8]{inputenc}
  \usepackage[T1]{fontenc}
  \usepackage{lmodern}
  \usepackage[portuguese]{babel}
  \DeclareUnicodeCharacter{2192}{\ensuremath{\rightarrow}}
\fi\fi
% raiz do repo: compilado da raiz ou de volN/<parte>/ (dois níveis abaixo)
\IfFileExists{tex/preambulo.tex}{% Preâmbulo comum ao livro e à pré-visualização de fragmentos (tex/_preview.tex).

% Mancha em A5.  Sem isto vale o padrão da classe book (~105 mm de largura e margem
% externa de ~25 mm em frente e verso), largo demais para o formato.  Aqui: ~118 mm de
% largura, margem interna um pouco maior que a externa (encadernação).
\usepackage[a5paper,inner=16mm,outer=14mm,top=14mm,bottom=16mm,
            includehead,headsep=5mm,footskip=8mm]{geometry}

% Títulos de secção sem justificar (evita linhas estouradas em \Large bold, p. ex.
% «Demonstrabilidade») e um pouco de folga para enunciados de tese em negrito.
\usepackage{sectsty}
\allsectionsfont{\raggedright}
\setlength{\emergencystretch}{1.5em}

% Marca da página impressa do original: \pg{286} imprime [286] no fluxo do texto.
\newcommand{\pg}[1]{\,{\footnotesize\textbf{[#1]}}\,}

% Virada de página do original DENTRO de uma nota (\pgn{422}); discreta, para não quebrar a nota.
\newcommand{\pgn}[1]{{\scriptsize\textbf{[#1]}}}

% Asterisco do original (divisão de assunto, sem título): \asterisco centrado em linha própria,
% também válido no meio do parágrafo ou dentro de nota (\footnote{... \asterisco ...}).
\newcommand{\asterisco}{\par\medskip\centerline{\footnotesize$\ast$}\medskip\noindent}

% Gravura recortada do fac-símile (<figure> do TEI; tools/tei2latex.py): fora de float, no ponto
% exato do original.  \gravura[detalhes]{arquivo.png}{fração da largura do texto}{legenda}:
% legenda principal centrada ("FIGURA VIII. — Côrte…"); os detalhes opcionais (<p> do <figure>:
% "A, hemispherio…; — B, …") vão abaixo, em corpo menor e justificados.  Ambos podem ser vazios.
\usepackage{graphicx}
\newcommand{\gravura}[4][]{\par\medskip\begingroup\centering
  \includegraphics[width=#3\linewidth,height=0.75\textheight,keepaspectratio]{#2}\par
  \if\relax\detokenize{#4}\relax\else\smallskip{\small #4\par}\fi
  \if\relax\detokenize{#1}\relax\else\smallskip{\footnotesize\leftskip=1em\rightskip=1em
    \parfillskip=0pt plus 1fil\noindent #1\par}\fi\endgroup\medskip\noindent}

% Citações latinas de S. Tomás (o OCR erra muito aqui: sempre conferir na imagem).
\newenvironment{latim}{\begin{quote}\itshape}{\end{quote}}

% --- Estrutura do autor (promovido de "Artigo III.tex") -----------------------
% enunciado de tese numerada que corre com o texto (segue-se travessão)
\newcommand{\tese}[2]{\bigskip\noindent\textbf{#1.\quad #2}\;\textemdash\;}
% enunciado de tese numerada em linha própria (seguem-se as subdivisões a), b)...)
\newcommand{\teseB}[2]{\bigskip\noindent\textbf{#1.\quad #2}\par}
% subdivisão a), b), c)... com rubrica em itálico
\newcommand{\sub}[2]{\medskip\noindent#1)\quad\emph{#2},\;\textemdash\;}
% subdivisão a), b), c)... sem rubrica
\newcommand{\subp}[1]{\medskip\noindent#1)\quad}

% --- Notas: no pé (padrão) ou correndo no corpo do texto ---------------------
% Passar  \def\notascorridas{}  no -usepretex faz cada nota correr no corpo, em
% corpo menor, num bloco recuado que interrompe a leitura até acabar.  Ex.:
%   latexmk -lualatex -usepretex='\def\fragmento{vol2/.../01_cap_y.tex}\def\notascorridas{}' tex/_preview.tex
\newif\ifnotascorridas
\ifdefined\notascorridas \notascorridastrue \fi
\ifnotascorridas
  \renewcommand\footnote[1]{%
    \stepcounter{footnote}\textsuperscript{\thefootnote}% marca no ponto do texto
    \par
    {\footnotesize
     \leftskip=2em \rightskip=2em \parindent=1em
     \addvspace{\medskipamount}%
     \noindent\textit{Nota \thefootnote.}\enspace #1\par
     \addvspace{\medskipamount}}%
  }
\fi

% --- Sem notas: \def\semnotas{} no -usepretex suprime todas as notas ---------
% (nem a chamada no texto nem o conteúdo).  Usado por  ./compila.sh --sem-notas .
\ifdefined\semnotas
  \renewcommand\footnote[2][]{}%
  \renewcommand\footnotemark[1][]{}%
  \renewcommand\footnotetext[2][]{}%
\fi

% peças pré-textuais (nó pretextual; p@type em web/obra/esquema.py)
\newcommand{\preTituloObra}[1]{\par\vspace{1em}{\centering\Large\bfseries\scshape #1\par}\vspace{1em}}
\newcommand{\preCentro}[1]{\par{\centering #1\par}}
\newcommand{\preEnderecamento}[1]{\par\vspace{1em}{\centering\small\scshape #1\par}}
\newcommand{\preSaudacao}[1]{\par\medskip\noindent #1\par}
\newcommand{\preDataLocal}[1]{\par\medskip\noindent\emph{#1}\par}
\newcommand{\preFecho}[1]{\par\smallskip{\raggedleft\scshape #1\par}}
}{% Preâmbulo comum ao livro e à pré-visualização de fragmentos (tex/_preview.tex).

% Mancha em A5.  Sem isto vale o padrão da classe book (~105 mm de largura e margem
% externa de ~25 mm em frente e verso), largo demais para o formato.  Aqui: ~118 mm de
% largura, margem interna um pouco maior que a externa (encadernação).
\usepackage[a5paper,inner=16mm,outer=14mm,top=14mm,bottom=16mm,
            includehead,headsep=5mm,footskip=8mm]{geometry}

% Títulos de secção sem justificar (evita linhas estouradas em \Large bold, p. ex.
% «Demonstrabilidade») e um pouco de folga para enunciados de tese em negrito.
\usepackage{sectsty}
\allsectionsfont{\raggedright}
\setlength{\emergencystretch}{1.5em}

% Marca da página impressa do original: \pg{286} imprime [286] no fluxo do texto.
\newcommand{\pg}[1]{\,{\footnotesize\textbf{[#1]}}\,}

% Virada de página do original DENTRO de uma nota (\pgn{422}); discreta, para não quebrar a nota.
\newcommand{\pgn}[1]{{\scriptsize\textbf{[#1]}}}

% Asterisco do original (divisão de assunto, sem título): \asterisco centrado em linha própria,
% também válido no meio do parágrafo ou dentro de nota (\footnote{... \asterisco ...}).
\newcommand{\asterisco}{\par\medskip\centerline{\footnotesize$\ast$}\medskip\noindent}

% Gravura recortada do fac-símile (<figure> do TEI; tools/tei2latex.py): fora de float, no ponto
% exato do original.  \gravura[detalhes]{arquivo.png}{fração da largura do texto}{legenda}:
% legenda principal centrada ("FIGURA VIII. — Côrte…"); os detalhes opcionais (<p> do <figure>:
% "A, hemispherio…; — B, …") vão abaixo, em corpo menor e justificados.  Ambos podem ser vazios.
\usepackage{graphicx}
\newcommand{\gravura}[4][]{\par\medskip\begingroup\centering
  \includegraphics[width=#3\linewidth,height=0.75\textheight,keepaspectratio]{#2}\par
  \if\relax\detokenize{#4}\relax\else\smallskip{\small #4\par}\fi
  \if\relax\detokenize{#1}\relax\else\smallskip{\footnotesize\leftskip=1em\rightskip=1em
    \parfillskip=0pt plus 1fil\noindent #1\par}\fi\endgroup\medskip\noindent}

% Citações latinas de S. Tomás (o OCR erra muito aqui: sempre conferir na imagem).
\newenvironment{latim}{\begin{quote}\itshape}{\end{quote}}

% --- Estrutura do autor (promovido de "Artigo III.tex") -----------------------
% enunciado de tese numerada que corre com o texto (segue-se travessão)
\newcommand{\tese}[2]{\bigskip\noindent\textbf{#1.\quad #2}\;\textemdash\;}
% enunciado de tese numerada em linha própria (seguem-se as subdivisões a), b)...)
\newcommand{\teseB}[2]{\bigskip\noindent\textbf{#1.\quad #2}\par}
% subdivisão a), b), c)... com rubrica em itálico
\newcommand{\sub}[2]{\medskip\noindent#1)\quad\emph{#2},\;\textemdash\;}
% subdivisão a), b), c)... sem rubrica
\newcommand{\subp}[1]{\medskip\noindent#1)\quad}

% --- Notas: no pé (padrão) ou correndo no corpo do texto ---------------------
% Passar  \def\notascorridas{}  no -usepretex faz cada nota correr no corpo, em
% corpo menor, num bloco recuado que interrompe a leitura até acabar.  Ex.:
%   latexmk -lualatex -usepretex='\def\fragmento{vol2/.../01_cap_y.tex}\def\notascorridas{}' tex/_preview.tex
\newif\ifnotascorridas
\ifdefined\notascorridas \notascorridastrue \fi
\ifnotascorridas
  \renewcommand\footnote[1]{%
    \stepcounter{footnote}\textsuperscript{\thefootnote}% marca no ponto do texto
    \par
    {\footnotesize
     \leftskip=2em \rightskip=2em \parindent=1em
     \addvspace{\medskipamount}%
     \noindent\textit{Nota \thefootnote.}\enspace #1\par
     \addvspace{\medskipamount}}%
  }
\fi

% --- Sem notas: \def\semnotas{} no -usepretex suprime todas as notas ---------
% (nem a chamada no texto nem o conteúdo).  Usado por  ./compila.sh --sem-notas .
\ifdefined\semnotas
  \renewcommand\footnote[2][]{}%
  \renewcommand\footnotemark[1][]{}%
  \renewcommand\footnotetext[2][]{}%
\fi

% peças pré-textuais (nó pretextual; p@type em web/obra/esquema.py)
\newcommand{\preTituloObra}[1]{\par\vspace{1em}{\centering\Large\bfseries\scshape #1\par}\vspace{1em}}
\newcommand{\preCentro}[1]{\par{\centering #1\par}}
\newcommand{\preEnderecamento}[1]{\par\vspace{1em}{\centering\small\scshape #1\par}}
\newcommand{\preSaudacao}[1]{\par\medskip\noindent #1\par}
\newcommand{\preDataLocal}[1]{\par\medskip\noindent\emph{#1}\par}
\newcommand{\preFecho}[1]{\par\smallskip{\raggedleft\scshape #1\par}}
}
\def\fimavulso{\end{document}}
\begin{document}
}{% Compilação AVULSA de um fragmento (volN/.../<stem>.tex) — p. ex. pelo botão do editor,
% de dentro da pasta do nó.  O fragmento começa com
%
%   \ifdefined\pg\else\IfFileExists{tex/avulso.tex}{% Compilação AVULSA de um fragmento (volN/.../<stem>.tex) — p. ex. pelo botão do editor,
% de dentro da pasta do nó.  O fragmento começa com
%
%   \ifdefined\pg\else\IfFileExists{tex/avulso.tex}{\input{tex/avulso}}{\input{../../tex/avulso}}\fi
%
% e termina com  \csname fimavulso\endcsname .  Dentro do livro ou do _preview, \pg já existe
% e nada disto roda; \fimavulso fica indefinido e o \csname vira \relax.
%
% Motor: o projeto é lualatex (fontspec/polyglossia).  Se o editor chamar pdflatex, cai
% num preâmbulo equivalente com inputenc/babel — serve para conferir, não para o PDF final.
\documentclass[11pt,a5paper]{book}
\ifdefined\directlua
  \usepackage{fontspec}
  \usepackage{polyglossia}
  \setmainlanguage{portuguese}
\else\ifdefined\XeTeXrevision
  \usepackage{fontspec}
  \usepackage{polyglossia}
  \setmainlanguage{portuguese}
\else
  \usepackage[utf8]{inputenc}
  \usepackage[T1]{fontenc}
  \usepackage{lmodern}
  \usepackage[portuguese]{babel}
  \DeclareUnicodeCharacter{2192}{\ensuremath{\rightarrow}}
\fi\fi
% raiz do repo: compilado da raiz ou de volN/<parte>/ (dois níveis abaixo)
\IfFileExists{tex/preambulo.tex}{\input{tex/preambulo}}{\input{../../tex/preambulo}}
\def\fimavulso{\end{document}}
\begin{document}
}{% Compilação AVULSA de um fragmento (volN/.../<stem>.tex) — p. ex. pelo botão do editor,
% de dentro da pasta do nó.  O fragmento começa com
%
%   \ifdefined\pg\else\IfFileExists{tex/avulso.tex}{\input{tex/avulso}}{\input{../../tex/avulso}}\fi
%
% e termina com  \csname fimavulso\endcsname .  Dentro do livro ou do _preview, \pg já existe
% e nada disto roda; \fimavulso fica indefinido e o \csname vira \relax.
%
% Motor: o projeto é lualatex (fontspec/polyglossia).  Se o editor chamar pdflatex, cai
% num preâmbulo equivalente com inputenc/babel — serve para conferir, não para o PDF final.
\documentclass[11pt,a5paper]{book}
\ifdefined\directlua
  \usepackage{fontspec}
  \usepackage{polyglossia}
  \setmainlanguage{portuguese}
\else\ifdefined\XeTeXrevision
  \usepackage{fontspec}
  \usepackage{polyglossia}
  \setmainlanguage{portuguese}
\else
  \usepackage[utf8]{inputenc}
  \usepackage[T1]{fontenc}
  \usepackage{lmodern}
  \usepackage[portuguese]{babel}
  \DeclareUnicodeCharacter{2192}{\ensuremath{\rightarrow}}
\fi\fi
% raiz do repo: compilado da raiz ou de volN/<parte>/ (dois níveis abaixo)
\IfFileExists{tex/preambulo.tex}{\input{tex/preambulo}}{\input{../../tex/preambulo}}
\def\fimavulso{\end{document}}
\begin{document}
}\fi
%
% e termina com  \csname fimavulso\endcsname .  Dentro do livro ou do _preview, \pg já existe
% e nada disto roda; \fimavulso fica indefinido e o \csname vira \relax.
%
% Motor: o projeto é lualatex (fontspec/polyglossia).  Se o editor chamar pdflatex, cai
% num preâmbulo equivalente com inputenc/babel — serve para conferir, não para o PDF final.
\documentclass[11pt,a5paper]{book}
\ifdefined\directlua
  \usepackage{fontspec}
  \usepackage{polyglossia}
  \setmainlanguage{portuguese}
\else\ifdefined\XeTeXrevision
  \usepackage{fontspec}
  \usepackage{polyglossia}
  \setmainlanguage{portuguese}
\else
  \usepackage[utf8]{inputenc}
  \usepackage[T1]{fontenc}
  \usepackage{lmodern}
  \usepackage[portuguese]{babel}
  \DeclareUnicodeCharacter{2192}{\ensuremath{\rightarrow}}
\fi\fi
% raiz do repo: compilado da raiz ou de volN/<parte>/ (dois níveis abaixo)
\IfFileExists{tex/preambulo.tex}{% Preâmbulo comum ao livro e à pré-visualização de fragmentos (tex/_preview.tex).

% Mancha em A5.  Sem isto vale o padrão da classe book (~105 mm de largura e margem
% externa de ~25 mm em frente e verso), largo demais para o formato.  Aqui: ~118 mm de
% largura, margem interna um pouco maior que a externa (encadernação).
\usepackage[a5paper,inner=16mm,outer=14mm,top=14mm,bottom=16mm,
            includehead,headsep=5mm,footskip=8mm]{geometry}

% Títulos de secção sem justificar (evita linhas estouradas em \Large bold, p. ex.
% «Demonstrabilidade») e um pouco de folga para enunciados de tese em negrito.
\usepackage{sectsty}
\allsectionsfont{\raggedright}
\setlength{\emergencystretch}{1.5em}

% Marca da página impressa do original: \pg{286} imprime [286] no fluxo do texto.
\newcommand{\pg}[1]{\,{\footnotesize\textbf{[#1]}}\,}

% Virada de página do original DENTRO de uma nota (\pgn{422}); discreta, para não quebrar a nota.
\newcommand{\pgn}[1]{{\scriptsize\textbf{[#1]}}}

% Asterisco do original (divisão de assunto, sem título): \asterisco centrado em linha própria,
% também válido no meio do parágrafo ou dentro de nota (\footnote{... \asterisco ...}).
\newcommand{\asterisco}{\par\medskip\centerline{\footnotesize$\ast$}\medskip\noindent}

% Gravura recortada do fac-símile (<figure> do TEI; tools/tei2latex.py): fora de float, no ponto
% exato do original.  \gravura[detalhes]{arquivo.png}{fração da largura do texto}{legenda}:
% legenda principal centrada ("FIGURA VIII. — Côrte…"); os detalhes opcionais (<p> do <figure>:
% "A, hemispherio…; — B, …") vão abaixo, em corpo menor e justificados.  Ambos podem ser vazios.
\usepackage{graphicx}
\newcommand{\gravura}[4][]{\par\medskip\begingroup\centering
  \includegraphics[width=#3\linewidth,height=0.75\textheight,keepaspectratio]{#2}\par
  \if\relax\detokenize{#4}\relax\else\smallskip{\small #4\par}\fi
  \if\relax\detokenize{#1}\relax\else\smallskip{\footnotesize\leftskip=1em\rightskip=1em
    \parfillskip=0pt plus 1fil\noindent #1\par}\fi\endgroup\medskip\noindent}

% Citações latinas de S. Tomás (o OCR erra muito aqui: sempre conferir na imagem).
\newenvironment{latim}{\begin{quote}\itshape}{\end{quote}}

% --- Estrutura do autor (promovido de "Artigo III.tex") -----------------------
% enunciado de tese numerada que corre com o texto (segue-se travessão)
\newcommand{\tese}[2]{\bigskip\noindent\textbf{#1.\quad #2}\;\textemdash\;}
% enunciado de tese numerada em linha própria (seguem-se as subdivisões a), b)...)
\newcommand{\teseB}[2]{\bigskip\noindent\textbf{#1.\quad #2}\par}
% subdivisão a), b), c)... com rubrica em itálico
\newcommand{\sub}[2]{\medskip\noindent#1)\quad\emph{#2},\;\textemdash\;}
% subdivisão a), b), c)... sem rubrica
\newcommand{\subp}[1]{\medskip\noindent#1)\quad}

% --- Notas: no pé (padrão) ou correndo no corpo do texto ---------------------
% Passar  \def\notascorridas{}  no -usepretex faz cada nota correr no corpo, em
% corpo menor, num bloco recuado que interrompe a leitura até acabar.  Ex.:
%   latexmk -lualatex -usepretex='\def\fragmento{vol2/.../01_cap_y.tex}\def\notascorridas{}' tex/_preview.tex
\newif\ifnotascorridas
\ifdefined\notascorridas \notascorridastrue \fi
\ifnotascorridas
  \renewcommand\footnote[1]{%
    \stepcounter{footnote}\textsuperscript{\thefootnote}% marca no ponto do texto
    \par
    {\footnotesize
     \leftskip=2em \rightskip=2em \parindent=1em
     \addvspace{\medskipamount}%
     \noindent\textit{Nota \thefootnote.}\enspace #1\par
     \addvspace{\medskipamount}}%
  }
\fi

% --- Sem notas: \def\semnotas{} no -usepretex suprime todas as notas ---------
% (nem a chamada no texto nem o conteúdo).  Usado por  ./compila.sh --sem-notas .
\ifdefined\semnotas
  \renewcommand\footnote[2][]{}%
  \renewcommand\footnotemark[1][]{}%
  \renewcommand\footnotetext[2][]{}%
\fi

% peças pré-textuais (nó pretextual; p@type em web/obra/esquema.py)
\newcommand{\preTituloObra}[1]{\par\vspace{1em}{\centering\Large\bfseries\scshape #1\par}\vspace{1em}}
\newcommand{\preCentro}[1]{\par{\centering #1\par}}
\newcommand{\preEnderecamento}[1]{\par\vspace{1em}{\centering\small\scshape #1\par}}
\newcommand{\preSaudacao}[1]{\par\medskip\noindent #1\par}
\newcommand{\preDataLocal}[1]{\par\medskip\noindent\emph{#1}\par}
\newcommand{\preFecho}[1]{\par\smallskip{\raggedleft\scshape #1\par}}
}{% Preâmbulo comum ao livro e à pré-visualização de fragmentos (tex/_preview.tex).

% Mancha em A5.  Sem isto vale o padrão da classe book (~105 mm de largura e margem
% externa de ~25 mm em frente e verso), largo demais para o formato.  Aqui: ~118 mm de
% largura, margem interna um pouco maior que a externa (encadernação).
\usepackage[a5paper,inner=16mm,outer=14mm,top=14mm,bottom=16mm,
            includehead,headsep=5mm,footskip=8mm]{geometry}

% Títulos de secção sem justificar (evita linhas estouradas em \Large bold, p. ex.
% «Demonstrabilidade») e um pouco de folga para enunciados de tese em negrito.
\usepackage{sectsty}
\allsectionsfont{\raggedright}
\setlength{\emergencystretch}{1.5em}

% Marca da página impressa do original: \pg{286} imprime [286] no fluxo do texto.
\newcommand{\pg}[1]{\,{\footnotesize\textbf{[#1]}}\,}

% Virada de página do original DENTRO de uma nota (\pgn{422}); discreta, para não quebrar a nota.
\newcommand{\pgn}[1]{{\scriptsize\textbf{[#1]}}}

% Asterisco do original (divisão de assunto, sem título): \asterisco centrado em linha própria,
% também válido no meio do parágrafo ou dentro de nota (\footnote{... \asterisco ...}).
\newcommand{\asterisco}{\par\medskip\centerline{\footnotesize$\ast$}\medskip\noindent}

% Gravura recortada do fac-símile (<figure> do TEI; tools/tei2latex.py): fora de float, no ponto
% exato do original.  \gravura[detalhes]{arquivo.png}{fração da largura do texto}{legenda}:
% legenda principal centrada ("FIGURA VIII. — Côrte…"); os detalhes opcionais (<p> do <figure>:
% "A, hemispherio…; — B, …") vão abaixo, em corpo menor e justificados.  Ambos podem ser vazios.
\usepackage{graphicx}
\newcommand{\gravura}[4][]{\par\medskip\begingroup\centering
  \includegraphics[width=#3\linewidth,height=0.75\textheight,keepaspectratio]{#2}\par
  \if\relax\detokenize{#4}\relax\else\smallskip{\small #4\par}\fi
  \if\relax\detokenize{#1}\relax\else\smallskip{\footnotesize\leftskip=1em\rightskip=1em
    \parfillskip=0pt plus 1fil\noindent #1\par}\fi\endgroup\medskip\noindent}

% Citações latinas de S. Tomás (o OCR erra muito aqui: sempre conferir na imagem).
\newenvironment{latim}{\begin{quote}\itshape}{\end{quote}}

% --- Estrutura do autor (promovido de "Artigo III.tex") -----------------------
% enunciado de tese numerada que corre com o texto (segue-se travessão)
\newcommand{\tese}[2]{\bigskip\noindent\textbf{#1.\quad #2}\;\textemdash\;}
% enunciado de tese numerada em linha própria (seguem-se as subdivisões a), b)...)
\newcommand{\teseB}[2]{\bigskip\noindent\textbf{#1.\quad #2}\par}
% subdivisão a), b), c)... com rubrica em itálico
\newcommand{\sub}[2]{\medskip\noindent#1)\quad\emph{#2},\;\textemdash\;}
% subdivisão a), b), c)... sem rubrica
\newcommand{\subp}[1]{\medskip\noindent#1)\quad}

% --- Notas: no pé (padrão) ou correndo no corpo do texto ---------------------
% Passar  \def\notascorridas{}  no -usepretex faz cada nota correr no corpo, em
% corpo menor, num bloco recuado que interrompe a leitura até acabar.  Ex.:
%   latexmk -lualatex -usepretex='\def\fragmento{vol2/.../01_cap_y.tex}\def\notascorridas{}' tex/_preview.tex
\newif\ifnotascorridas
\ifdefined\notascorridas \notascorridastrue \fi
\ifnotascorridas
  \renewcommand\footnote[1]{%
    \stepcounter{footnote}\textsuperscript{\thefootnote}% marca no ponto do texto
    \par
    {\footnotesize
     \leftskip=2em \rightskip=2em \parindent=1em
     \addvspace{\medskipamount}%
     \noindent\textit{Nota \thefootnote.}\enspace #1\par
     \addvspace{\medskipamount}}%
  }
\fi

% --- Sem notas: \def\semnotas{} no -usepretex suprime todas as notas ---------
% (nem a chamada no texto nem o conteúdo).  Usado por  ./compila.sh --sem-notas .
\ifdefined\semnotas
  \renewcommand\footnote[2][]{}%
  \renewcommand\footnotemark[1][]{}%
  \renewcommand\footnotetext[2][]{}%
\fi

% peças pré-textuais (nó pretextual; p@type em web/obra/esquema.py)
\newcommand{\preTituloObra}[1]{\par\vspace{1em}{\centering\Large\bfseries\scshape #1\par}\vspace{1em}}
\newcommand{\preCentro}[1]{\par{\centering #1\par}}
\newcommand{\preEnderecamento}[1]{\par\vspace{1em}{\centering\small\scshape #1\par}}
\newcommand{\preSaudacao}[1]{\par\medskip\noindent #1\par}
\newcommand{\preDataLocal}[1]{\par\medskip\noindent\emph{#1}\par}
\newcommand{\preFecho}[1]{\par\smallskip{\raggedleft\scshape #1\par}}
}
\def\fimavulso{\end{document}}
\begin{document}
}\fi
%
% e termina com  \csname fimavulso\endcsname .  Dentro do livro ou do _preview, \pg já existe
% e nada disto roda; \fimavulso fica indefinido e o \csname vira \relax.
%
% Motor: o projeto é lualatex (fontspec/polyglossia).  Se o editor chamar pdflatex, cai
% num preâmbulo equivalente com inputenc/babel — serve para conferir, não para o PDF final.
\documentclass[11pt,a5paper]{book}
\ifdefined\directlua
  \usepackage{fontspec}
  \usepackage{polyglossia}
  \setmainlanguage{portuguese}
\else\ifdefined\XeTeXrevision
  \usepackage{fontspec}
  \usepackage{polyglossia}
  \setmainlanguage{portuguese}
\else
  \usepackage[utf8]{inputenc}
  \usepackage[T1]{fontenc}
  \usepackage{lmodern}
  \usepackage[portuguese]{babel}
  \DeclareUnicodeCharacter{2192}{\ensuremath{\rightarrow}}
\fi\fi
% raiz do repo: compilado da raiz ou de volN/<parte>/ (dois níveis abaixo)
\IfFileExists{tex/preambulo.tex}{% Preâmbulo comum ao livro e à pré-visualização de fragmentos (tex/_preview.tex).

% Mancha em A5.  Sem isto vale o padrão da classe book (~105 mm de largura e margem
% externa de ~25 mm em frente e verso), largo demais para o formato.  Aqui: ~118 mm de
% largura, margem interna um pouco maior que a externa (encadernação).
\usepackage[a5paper,inner=16mm,outer=14mm,top=14mm,bottom=16mm,
            includehead,headsep=5mm,footskip=8mm]{geometry}

% Títulos de secção sem justificar (evita linhas estouradas em \Large bold, p. ex.
% «Demonstrabilidade») e um pouco de folga para enunciados de tese em negrito.
\usepackage{sectsty}
\allsectionsfont{\raggedright}
\setlength{\emergencystretch}{1.5em}

% Marca da página impressa do original: \pg{286} imprime [286] no fluxo do texto.
\newcommand{\pg}[1]{\,{\footnotesize\textbf{[#1]}}\,}

% Virada de página do original DENTRO de uma nota (\pgn{422}); discreta, para não quebrar a nota.
\newcommand{\pgn}[1]{{\scriptsize\textbf{[#1]}}}

% Asterisco do original (divisão de assunto, sem título): \asterisco centrado em linha própria,
% também válido no meio do parágrafo ou dentro de nota (\footnote{... \asterisco ...}).
\newcommand{\asterisco}{\par\medskip\centerline{\footnotesize$\ast$}\medskip\noindent}

% Gravura recortada do fac-símile (<figure> do TEI; tools/tei2latex.py): fora de float, no ponto
% exato do original.  \gravura[detalhes]{arquivo.png}{fração da largura do texto}{legenda}:
% legenda principal centrada ("FIGURA VIII. — Côrte…"); os detalhes opcionais (<p> do <figure>:
% "A, hemispherio…; — B, …") vão abaixo, em corpo menor e justificados.  Ambos podem ser vazios.
\usepackage{graphicx}
\newcommand{\gravura}[4][]{\par\medskip\begingroup\centering
  \includegraphics[width=#3\linewidth,height=0.75\textheight,keepaspectratio]{#2}\par
  \if\relax\detokenize{#4}\relax\else\smallskip{\small #4\par}\fi
  \if\relax\detokenize{#1}\relax\else\smallskip{\footnotesize\leftskip=1em\rightskip=1em
    \parfillskip=0pt plus 1fil\noindent #1\par}\fi\endgroup\medskip\noindent}

% Citações latinas de S. Tomás (o OCR erra muito aqui: sempre conferir na imagem).
\newenvironment{latim}{\begin{quote}\itshape}{\end{quote}}

% --- Estrutura do autor (promovido de "Artigo III.tex") -----------------------
% enunciado de tese numerada que corre com o texto (segue-se travessão)
\newcommand{\tese}[2]{\bigskip\noindent\textbf{#1.\quad #2}\;\textemdash\;}
% enunciado de tese numerada em linha própria (seguem-se as subdivisões a), b)...)
\newcommand{\teseB}[2]{\bigskip\noindent\textbf{#1.\quad #2}\par}
% subdivisão a), b), c)... com rubrica em itálico
\newcommand{\sub}[2]{\medskip\noindent#1)\quad\emph{#2},\;\textemdash\;}
% subdivisão a), b), c)... sem rubrica
\newcommand{\subp}[1]{\medskip\noindent#1)\quad}

% --- Notas: no pé (padrão) ou correndo no corpo do texto ---------------------
% Passar  \def\notascorridas{}  no -usepretex faz cada nota correr no corpo, em
% corpo menor, num bloco recuado que interrompe a leitura até acabar.  Ex.:
%   latexmk -lualatex -usepretex='\def\fragmento{vol2/.../01_cap_y.tex}\def\notascorridas{}' tex/_preview.tex
\newif\ifnotascorridas
\ifdefined\notascorridas \notascorridastrue \fi
\ifnotascorridas
  \renewcommand\footnote[1]{%
    \stepcounter{footnote}\textsuperscript{\thefootnote}% marca no ponto do texto
    \par
    {\footnotesize
     \leftskip=2em \rightskip=2em \parindent=1em
     \addvspace{\medskipamount}%
     \noindent\textit{Nota \thefootnote.}\enspace #1\par
     \addvspace{\medskipamount}}%
  }
\fi

% --- Sem notas: \def\semnotas{} no -usepretex suprime todas as notas ---------
% (nem a chamada no texto nem o conteúdo).  Usado por  ./compila.sh --sem-notas .
\ifdefined\semnotas
  \renewcommand\footnote[2][]{}%
  \renewcommand\footnotemark[1][]{}%
  \renewcommand\footnotetext[2][]{}%
\fi

% peças pré-textuais (nó pretextual; p@type em web/obra/esquema.py)
\newcommand{\preTituloObra}[1]{\par\vspace{1em}{\centering\Large\bfseries\scshape #1\par}\vspace{1em}}
\newcommand{\preCentro}[1]{\par{\centering #1\par}}
\newcommand{\preEnderecamento}[1]{\par\vspace{1em}{\centering\small\scshape #1\par}}
\newcommand{\preSaudacao}[1]{\par\medskip\noindent #1\par}
\newcommand{\preDataLocal}[1]{\par\medskip\noindent\emph{#1}\par}
\newcommand{\preFecho}[1]{\par\smallskip{\raggedleft\scshape #1\par}}
}{% Preâmbulo comum ao livro e à pré-visualização de fragmentos (tex/_preview.tex).

% Mancha em A5.  Sem isto vale o padrão da classe book (~105 mm de largura e margem
% externa de ~25 mm em frente e verso), largo demais para o formato.  Aqui: ~118 mm de
% largura, margem interna um pouco maior que a externa (encadernação).
\usepackage[a5paper,inner=16mm,outer=14mm,top=14mm,bottom=16mm,
            includehead,headsep=5mm,footskip=8mm]{geometry}

% Títulos de secção sem justificar (evita linhas estouradas em \Large bold, p. ex.
% «Demonstrabilidade») e um pouco de folga para enunciados de tese em negrito.
\usepackage{sectsty}
\allsectionsfont{\raggedright}
\setlength{\emergencystretch}{1.5em}

% Marca da página impressa do original: \pg{286} imprime [286] no fluxo do texto.
\newcommand{\pg}[1]{\,{\footnotesize\textbf{[#1]}}\,}

% Virada de página do original DENTRO de uma nota (\pgn{422}); discreta, para não quebrar a nota.
\newcommand{\pgn}[1]{{\scriptsize\textbf{[#1]}}}

% Asterisco do original (divisão de assunto, sem título): \asterisco centrado em linha própria,
% também válido no meio do parágrafo ou dentro de nota (\footnote{... \asterisco ...}).
\newcommand{\asterisco}{\par\medskip\centerline{\footnotesize$\ast$}\medskip\noindent}

% Gravura recortada do fac-símile (<figure> do TEI; tools/tei2latex.py): fora de float, no ponto
% exato do original.  \gravura[detalhes]{arquivo.png}{fração da largura do texto}{legenda}:
% legenda principal centrada ("FIGURA VIII. — Côrte…"); os detalhes opcionais (<p> do <figure>:
% "A, hemispherio…; — B, …") vão abaixo, em corpo menor e justificados.  Ambos podem ser vazios.
\usepackage{graphicx}
\newcommand{\gravura}[4][]{\par\medskip\begingroup\centering
  \includegraphics[width=#3\linewidth,height=0.75\textheight,keepaspectratio]{#2}\par
  \if\relax\detokenize{#4}\relax\else\smallskip{\small #4\par}\fi
  \if\relax\detokenize{#1}\relax\else\smallskip{\footnotesize\leftskip=1em\rightskip=1em
    \parfillskip=0pt plus 1fil\noindent #1\par}\fi\endgroup\medskip\noindent}

% Citações latinas de S. Tomás (o OCR erra muito aqui: sempre conferir na imagem).
\newenvironment{latim}{\begin{quote}\itshape}{\end{quote}}

% --- Estrutura do autor (promovido de "Artigo III.tex") -----------------------
% enunciado de tese numerada que corre com o texto (segue-se travessão)
\newcommand{\tese}[2]{\bigskip\noindent\textbf{#1.\quad #2}\;\textemdash\;}
% enunciado de tese numerada em linha própria (seguem-se as subdivisões a), b)...)
\newcommand{\teseB}[2]{\bigskip\noindent\textbf{#1.\quad #2}\par}
% subdivisão a), b), c)... com rubrica em itálico
\newcommand{\sub}[2]{\medskip\noindent#1)\quad\emph{#2},\;\textemdash\;}
% subdivisão a), b), c)... sem rubrica
\newcommand{\subp}[1]{\medskip\noindent#1)\quad}

% --- Notas: no pé (padrão) ou correndo no corpo do texto ---------------------
% Passar  \def\notascorridas{}  no -usepretex faz cada nota correr no corpo, em
% corpo menor, num bloco recuado que interrompe a leitura até acabar.  Ex.:
%   latexmk -lualatex -usepretex='\def\fragmento{vol2/.../01_cap_y.tex}\def\notascorridas{}' tex/_preview.tex
\newif\ifnotascorridas
\ifdefined\notascorridas \notascorridastrue \fi
\ifnotascorridas
  \renewcommand\footnote[1]{%
    \stepcounter{footnote}\textsuperscript{\thefootnote}% marca no ponto do texto
    \par
    {\footnotesize
     \leftskip=2em \rightskip=2em \parindent=1em
     \addvspace{\medskipamount}%
     \noindent\textit{Nota \thefootnote.}\enspace #1\par
     \addvspace{\medskipamount}}%
  }
\fi

% --- Sem notas: \def\semnotas{} no -usepretex suprime todas as notas ---------
% (nem a chamada no texto nem o conteúdo).  Usado por  ./compila.sh --sem-notas .
\ifdefined\semnotas
  \renewcommand\footnote[2][]{}%
  \renewcommand\footnotemark[1][]{}%
  \renewcommand\footnotetext[2][]{}%
\fi

% peças pré-textuais (nó pretextual; p@type em web/obra/esquema.py)
\newcommand{\preTituloObra}[1]{\par\vspace{1em}{\centering\Large\bfseries\scshape #1\par}\vspace{1em}}
\newcommand{\preCentro}[1]{\par{\centering #1\par}}
\newcommand{\preEnderecamento}[1]{\par\vspace{1em}{\centering\small\scshape #1\par}}
\newcommand{\preSaudacao}[1]{\par\medskip\noindent #1\par}
\newcommand{\preDataLocal}[1]{\par\medskip\noindent\emph{#1}\par}
\newcommand{\preFecho}[1]{\par\smallskip{\raggedleft\scshape #1\par}}
}
\def\fimavulso{\end{document}}
\begin{document}
}{% Compilação AVULSA de um fragmento (volN/.../<stem>.tex) — p. ex. pelo botão do editor,
% de dentro da pasta do nó.  O fragmento começa com
%
%   \ifdefined\pg\else\IfFileExists{tex/avulso.tex}{% Compilação AVULSA de um fragmento (volN/.../<stem>.tex) — p. ex. pelo botão do editor,
% de dentro da pasta do nó.  O fragmento começa com
%
%   \ifdefined\pg\else\IfFileExists{tex/avulso.tex}{% Compilação AVULSA de um fragmento (volN/.../<stem>.tex) — p. ex. pelo botão do editor,
% de dentro da pasta do nó.  O fragmento começa com
%
%   \ifdefined\pg\else\IfFileExists{tex/avulso.tex}{\input{tex/avulso}}{\input{../../tex/avulso}}\fi
%
% e termina com  \csname fimavulso\endcsname .  Dentro do livro ou do _preview, \pg já existe
% e nada disto roda; \fimavulso fica indefinido e o \csname vira \relax.
%
% Motor: o projeto é lualatex (fontspec/polyglossia).  Se o editor chamar pdflatex, cai
% num preâmbulo equivalente com inputenc/babel — serve para conferir, não para o PDF final.
\documentclass[11pt,a5paper]{book}
\ifdefined\directlua
  \usepackage{fontspec}
  \usepackage{polyglossia}
  \setmainlanguage{portuguese}
\else\ifdefined\XeTeXrevision
  \usepackage{fontspec}
  \usepackage{polyglossia}
  \setmainlanguage{portuguese}
\else
  \usepackage[utf8]{inputenc}
  \usepackage[T1]{fontenc}
  \usepackage{lmodern}
  \usepackage[portuguese]{babel}
  \DeclareUnicodeCharacter{2192}{\ensuremath{\rightarrow}}
\fi\fi
% raiz do repo: compilado da raiz ou de volN/<parte>/ (dois níveis abaixo)
\IfFileExists{tex/preambulo.tex}{\input{tex/preambulo}}{\input{../../tex/preambulo}}
\def\fimavulso{\end{document}}
\begin{document}
}{% Compilação AVULSA de um fragmento (volN/.../<stem>.tex) — p. ex. pelo botão do editor,
% de dentro da pasta do nó.  O fragmento começa com
%
%   \ifdefined\pg\else\IfFileExists{tex/avulso.tex}{\input{tex/avulso}}{\input{../../tex/avulso}}\fi
%
% e termina com  \csname fimavulso\endcsname .  Dentro do livro ou do _preview, \pg já existe
% e nada disto roda; \fimavulso fica indefinido e o \csname vira \relax.
%
% Motor: o projeto é lualatex (fontspec/polyglossia).  Se o editor chamar pdflatex, cai
% num preâmbulo equivalente com inputenc/babel — serve para conferir, não para o PDF final.
\documentclass[11pt,a5paper]{book}
\ifdefined\directlua
  \usepackage{fontspec}
  \usepackage{polyglossia}
  \setmainlanguage{portuguese}
\else\ifdefined\XeTeXrevision
  \usepackage{fontspec}
  \usepackage{polyglossia}
  \setmainlanguage{portuguese}
\else
  \usepackage[utf8]{inputenc}
  \usepackage[T1]{fontenc}
  \usepackage{lmodern}
  \usepackage[portuguese]{babel}
  \DeclareUnicodeCharacter{2192}{\ensuremath{\rightarrow}}
\fi\fi
% raiz do repo: compilado da raiz ou de volN/<parte>/ (dois níveis abaixo)
\IfFileExists{tex/preambulo.tex}{\input{tex/preambulo}}{\input{../../tex/preambulo}}
\def\fimavulso{\end{document}}
\begin{document}
}\fi
%
% e termina com  \csname fimavulso\endcsname .  Dentro do livro ou do _preview, \pg já existe
% e nada disto roda; \fimavulso fica indefinido e o \csname vira \relax.
%
% Motor: o projeto é lualatex (fontspec/polyglossia).  Se o editor chamar pdflatex, cai
% num preâmbulo equivalente com inputenc/babel — serve para conferir, não para o PDF final.
\documentclass[11pt,a5paper]{book}
\ifdefined\directlua
  \usepackage{fontspec}
  \usepackage{polyglossia}
  \setmainlanguage{portuguese}
\else\ifdefined\XeTeXrevision
  \usepackage{fontspec}
  \usepackage{polyglossia}
  \setmainlanguage{portuguese}
\else
  \usepackage[utf8]{inputenc}
  \usepackage[T1]{fontenc}
  \usepackage{lmodern}
  \usepackage[portuguese]{babel}
  \DeclareUnicodeCharacter{2192}{\ensuremath{\rightarrow}}
\fi\fi
% raiz do repo: compilado da raiz ou de volN/<parte>/ (dois níveis abaixo)
\IfFileExists{tex/preambulo.tex}{% Preâmbulo comum ao livro e à pré-visualização de fragmentos (tex/_preview.tex).

% Mancha em A5.  Sem isto vale o padrão da classe book (~105 mm de largura e margem
% externa de ~25 mm em frente e verso), largo demais para o formato.  Aqui: ~118 mm de
% largura, margem interna um pouco maior que a externa (encadernação).
\usepackage[a5paper,inner=16mm,outer=14mm,top=14mm,bottom=16mm,
            includehead,headsep=5mm,footskip=8mm]{geometry}

% Títulos de secção sem justificar (evita linhas estouradas em \Large bold, p. ex.
% «Demonstrabilidade») e um pouco de folga para enunciados de tese em negrito.
\usepackage{sectsty}
\allsectionsfont{\raggedright}
\setlength{\emergencystretch}{1.5em}

% Marca da página impressa do original: \pg{286} imprime [286] no fluxo do texto.
\newcommand{\pg}[1]{\,{\footnotesize\textbf{[#1]}}\,}

% Virada de página do original DENTRO de uma nota (\pgn{422}); discreta, para não quebrar a nota.
\newcommand{\pgn}[1]{{\scriptsize\textbf{[#1]}}}

% Asterisco do original (divisão de assunto, sem título): \asterisco centrado em linha própria,
% também válido no meio do parágrafo ou dentro de nota (\footnote{... \asterisco ...}).
\newcommand{\asterisco}{\par\medskip\centerline{\footnotesize$\ast$}\medskip\noindent}

% Gravura recortada do fac-símile (<figure> do TEI; tools/tei2latex.py): fora de float, no ponto
% exato do original.  \gravura[detalhes]{arquivo.png}{fração da largura do texto}{legenda}:
% legenda principal centrada ("FIGURA VIII. — Côrte…"); os detalhes opcionais (<p> do <figure>:
% "A, hemispherio…; — B, …") vão abaixo, em corpo menor e justificados.  Ambos podem ser vazios.
\usepackage{graphicx}
\newcommand{\gravura}[4][]{\par\medskip\begingroup\centering
  \includegraphics[width=#3\linewidth,height=0.75\textheight,keepaspectratio]{#2}\par
  \if\relax\detokenize{#4}\relax\else\smallskip{\small #4\par}\fi
  \if\relax\detokenize{#1}\relax\else\smallskip{\footnotesize\leftskip=1em\rightskip=1em
    \parfillskip=0pt plus 1fil\noindent #1\par}\fi\endgroup\medskip\noindent}

% Citações latinas de S. Tomás (o OCR erra muito aqui: sempre conferir na imagem).
\newenvironment{latim}{\begin{quote}\itshape}{\end{quote}}

% --- Estrutura do autor (promovido de "Artigo III.tex") -----------------------
% enunciado de tese numerada que corre com o texto (segue-se travessão)
\newcommand{\tese}[2]{\bigskip\noindent\textbf{#1.\quad #2}\;\textemdash\;}
% enunciado de tese numerada em linha própria (seguem-se as subdivisões a), b)...)
\newcommand{\teseB}[2]{\bigskip\noindent\textbf{#1.\quad #2}\par}
% subdivisão a), b), c)... com rubrica em itálico
\newcommand{\sub}[2]{\medskip\noindent#1)\quad\emph{#2},\;\textemdash\;}
% subdivisão a), b), c)... sem rubrica
\newcommand{\subp}[1]{\medskip\noindent#1)\quad}

% --- Notas: no pé (padrão) ou correndo no corpo do texto ---------------------
% Passar  \def\notascorridas{}  no -usepretex faz cada nota correr no corpo, em
% corpo menor, num bloco recuado que interrompe a leitura até acabar.  Ex.:
%   latexmk -lualatex -usepretex='\def\fragmento{vol2/.../01_cap_y.tex}\def\notascorridas{}' tex/_preview.tex
\newif\ifnotascorridas
\ifdefined\notascorridas \notascorridastrue \fi
\ifnotascorridas
  \renewcommand\footnote[1]{%
    \stepcounter{footnote}\textsuperscript{\thefootnote}% marca no ponto do texto
    \par
    {\footnotesize
     \leftskip=2em \rightskip=2em \parindent=1em
     \addvspace{\medskipamount}%
     \noindent\textit{Nota \thefootnote.}\enspace #1\par
     \addvspace{\medskipamount}}%
  }
\fi

% --- Sem notas: \def\semnotas{} no -usepretex suprime todas as notas ---------
% (nem a chamada no texto nem o conteúdo).  Usado por  ./compila.sh --sem-notas .
\ifdefined\semnotas
  \renewcommand\footnote[2][]{}%
  \renewcommand\footnotemark[1][]{}%
  \renewcommand\footnotetext[2][]{}%
\fi

% peças pré-textuais (nó pretextual; p@type em web/obra/esquema.py)
\newcommand{\preTituloObra}[1]{\par\vspace{1em}{\centering\Large\bfseries\scshape #1\par}\vspace{1em}}
\newcommand{\preCentro}[1]{\par{\centering #1\par}}
\newcommand{\preEnderecamento}[1]{\par\vspace{1em}{\centering\small\scshape #1\par}}
\newcommand{\preSaudacao}[1]{\par\medskip\noindent #1\par}
\newcommand{\preDataLocal}[1]{\par\medskip\noindent\emph{#1}\par}
\newcommand{\preFecho}[1]{\par\smallskip{\raggedleft\scshape #1\par}}
}{% Preâmbulo comum ao livro e à pré-visualização de fragmentos (tex/_preview.tex).

% Mancha em A5.  Sem isto vale o padrão da classe book (~105 mm de largura e margem
% externa de ~25 mm em frente e verso), largo demais para o formato.  Aqui: ~118 mm de
% largura, margem interna um pouco maior que a externa (encadernação).
\usepackage[a5paper,inner=16mm,outer=14mm,top=14mm,bottom=16mm,
            includehead,headsep=5mm,footskip=8mm]{geometry}

% Títulos de secção sem justificar (evita linhas estouradas em \Large bold, p. ex.
% «Demonstrabilidade») e um pouco de folga para enunciados de tese em negrito.
\usepackage{sectsty}
\allsectionsfont{\raggedright}
\setlength{\emergencystretch}{1.5em}

% Marca da página impressa do original: \pg{286} imprime [286] no fluxo do texto.
\newcommand{\pg}[1]{\,{\footnotesize\textbf{[#1]}}\,}

% Virada de página do original DENTRO de uma nota (\pgn{422}); discreta, para não quebrar a nota.
\newcommand{\pgn}[1]{{\scriptsize\textbf{[#1]}}}

% Asterisco do original (divisão de assunto, sem título): \asterisco centrado em linha própria,
% também válido no meio do parágrafo ou dentro de nota (\footnote{... \asterisco ...}).
\newcommand{\asterisco}{\par\medskip\centerline{\footnotesize$\ast$}\medskip\noindent}

% Gravura recortada do fac-símile (<figure> do TEI; tools/tei2latex.py): fora de float, no ponto
% exato do original.  \gravura[detalhes]{arquivo.png}{fração da largura do texto}{legenda}:
% legenda principal centrada ("FIGURA VIII. — Côrte…"); os detalhes opcionais (<p> do <figure>:
% "A, hemispherio…; — B, …") vão abaixo, em corpo menor e justificados.  Ambos podem ser vazios.
\usepackage{graphicx}
\newcommand{\gravura}[4][]{\par\medskip\begingroup\centering
  \includegraphics[width=#3\linewidth,height=0.75\textheight,keepaspectratio]{#2}\par
  \if\relax\detokenize{#4}\relax\else\smallskip{\small #4\par}\fi
  \if\relax\detokenize{#1}\relax\else\smallskip{\footnotesize\leftskip=1em\rightskip=1em
    \parfillskip=0pt plus 1fil\noindent #1\par}\fi\endgroup\medskip\noindent}

% Citações latinas de S. Tomás (o OCR erra muito aqui: sempre conferir na imagem).
\newenvironment{latim}{\begin{quote}\itshape}{\end{quote}}

% --- Estrutura do autor (promovido de "Artigo III.tex") -----------------------
% enunciado de tese numerada que corre com o texto (segue-se travessão)
\newcommand{\tese}[2]{\bigskip\noindent\textbf{#1.\quad #2}\;\textemdash\;}
% enunciado de tese numerada em linha própria (seguem-se as subdivisões a), b)...)
\newcommand{\teseB}[2]{\bigskip\noindent\textbf{#1.\quad #2}\par}
% subdivisão a), b), c)... com rubrica em itálico
\newcommand{\sub}[2]{\medskip\noindent#1)\quad\emph{#2},\;\textemdash\;}
% subdivisão a), b), c)... sem rubrica
\newcommand{\subp}[1]{\medskip\noindent#1)\quad}

% --- Notas: no pé (padrão) ou correndo no corpo do texto ---------------------
% Passar  \def\notascorridas{}  no -usepretex faz cada nota correr no corpo, em
% corpo menor, num bloco recuado que interrompe a leitura até acabar.  Ex.:
%   latexmk -lualatex -usepretex='\def\fragmento{vol2/.../01_cap_y.tex}\def\notascorridas{}' tex/_preview.tex
\newif\ifnotascorridas
\ifdefined\notascorridas \notascorridastrue \fi
\ifnotascorridas
  \renewcommand\footnote[1]{%
    \stepcounter{footnote}\textsuperscript{\thefootnote}% marca no ponto do texto
    \par
    {\footnotesize
     \leftskip=2em \rightskip=2em \parindent=1em
     \addvspace{\medskipamount}%
     \noindent\textit{Nota \thefootnote.}\enspace #1\par
     \addvspace{\medskipamount}}%
  }
\fi

% --- Sem notas: \def\semnotas{} no -usepretex suprime todas as notas ---------
% (nem a chamada no texto nem o conteúdo).  Usado por  ./compila.sh --sem-notas .
\ifdefined\semnotas
  \renewcommand\footnote[2][]{}%
  \renewcommand\footnotemark[1][]{}%
  \renewcommand\footnotetext[2][]{}%
\fi

% peças pré-textuais (nó pretextual; p@type em web/obra/esquema.py)
\newcommand{\preTituloObra}[1]{\par\vspace{1em}{\centering\Large\bfseries\scshape #1\par}\vspace{1em}}
\newcommand{\preCentro}[1]{\par{\centering #1\par}}
\newcommand{\preEnderecamento}[1]{\par\vspace{1em}{\centering\small\scshape #1\par}}
\newcommand{\preSaudacao}[1]{\par\medskip\noindent #1\par}
\newcommand{\preDataLocal}[1]{\par\medskip\noindent\emph{#1}\par}
\newcommand{\preFecho}[1]{\par\smallskip{\raggedleft\scshape #1\par}}
}
\def\fimavulso{\end{document}}
\begin{document}
}{% Compilação AVULSA de um fragmento (volN/.../<stem>.tex) — p. ex. pelo botão do editor,
% de dentro da pasta do nó.  O fragmento começa com
%
%   \ifdefined\pg\else\IfFileExists{tex/avulso.tex}{% Compilação AVULSA de um fragmento (volN/.../<stem>.tex) — p. ex. pelo botão do editor,
% de dentro da pasta do nó.  O fragmento começa com
%
%   \ifdefined\pg\else\IfFileExists{tex/avulso.tex}{\input{tex/avulso}}{\input{../../tex/avulso}}\fi
%
% e termina com  \csname fimavulso\endcsname .  Dentro do livro ou do _preview, \pg já existe
% e nada disto roda; \fimavulso fica indefinido e o \csname vira \relax.
%
% Motor: o projeto é lualatex (fontspec/polyglossia).  Se o editor chamar pdflatex, cai
% num preâmbulo equivalente com inputenc/babel — serve para conferir, não para o PDF final.
\documentclass[11pt,a5paper]{book}
\ifdefined\directlua
  \usepackage{fontspec}
  \usepackage{polyglossia}
  \setmainlanguage{portuguese}
\else\ifdefined\XeTeXrevision
  \usepackage{fontspec}
  \usepackage{polyglossia}
  \setmainlanguage{portuguese}
\else
  \usepackage[utf8]{inputenc}
  \usepackage[T1]{fontenc}
  \usepackage{lmodern}
  \usepackage[portuguese]{babel}
  \DeclareUnicodeCharacter{2192}{\ensuremath{\rightarrow}}
\fi\fi
% raiz do repo: compilado da raiz ou de volN/<parte>/ (dois níveis abaixo)
\IfFileExists{tex/preambulo.tex}{\input{tex/preambulo}}{\input{../../tex/preambulo}}
\def\fimavulso{\end{document}}
\begin{document}
}{% Compilação AVULSA de um fragmento (volN/.../<stem>.tex) — p. ex. pelo botão do editor,
% de dentro da pasta do nó.  O fragmento começa com
%
%   \ifdefined\pg\else\IfFileExists{tex/avulso.tex}{\input{tex/avulso}}{\input{../../tex/avulso}}\fi
%
% e termina com  \csname fimavulso\endcsname .  Dentro do livro ou do _preview, \pg já existe
% e nada disto roda; \fimavulso fica indefinido e o \csname vira \relax.
%
% Motor: o projeto é lualatex (fontspec/polyglossia).  Se o editor chamar pdflatex, cai
% num preâmbulo equivalente com inputenc/babel — serve para conferir, não para o PDF final.
\documentclass[11pt,a5paper]{book}
\ifdefined\directlua
  \usepackage{fontspec}
  \usepackage{polyglossia}
  \setmainlanguage{portuguese}
\else\ifdefined\XeTeXrevision
  \usepackage{fontspec}
  \usepackage{polyglossia}
  \setmainlanguage{portuguese}
\else
  \usepackage[utf8]{inputenc}
  \usepackage[T1]{fontenc}
  \usepackage{lmodern}
  \usepackage[portuguese]{babel}
  \DeclareUnicodeCharacter{2192}{\ensuremath{\rightarrow}}
\fi\fi
% raiz do repo: compilado da raiz ou de volN/<parte>/ (dois níveis abaixo)
\IfFileExists{tex/preambulo.tex}{\input{tex/preambulo}}{\input{../../tex/preambulo}}
\def\fimavulso{\end{document}}
\begin{document}
}\fi
%
% e termina com  \csname fimavulso\endcsname .  Dentro do livro ou do _preview, \pg já existe
% e nada disto roda; \fimavulso fica indefinido e o \csname vira \relax.
%
% Motor: o projeto é lualatex (fontspec/polyglossia).  Se o editor chamar pdflatex, cai
% num preâmbulo equivalente com inputenc/babel — serve para conferir, não para o PDF final.
\documentclass[11pt,a5paper]{book}
\ifdefined\directlua
  \usepackage{fontspec}
  \usepackage{polyglossia}
  \setmainlanguage{portuguese}
\else\ifdefined\XeTeXrevision
  \usepackage{fontspec}
  \usepackage{polyglossia}
  \setmainlanguage{portuguese}
\else
  \usepackage[utf8]{inputenc}
  \usepackage[T1]{fontenc}
  \usepackage{lmodern}
  \usepackage[portuguese]{babel}
  \DeclareUnicodeCharacter{2192}{\ensuremath{\rightarrow}}
\fi\fi
% raiz do repo: compilado da raiz ou de volN/<parte>/ (dois níveis abaixo)
\IfFileExists{tex/preambulo.tex}{% Preâmbulo comum ao livro e à pré-visualização de fragmentos (tex/_preview.tex).

% Mancha em A5.  Sem isto vale o padrão da classe book (~105 mm de largura e margem
% externa de ~25 mm em frente e verso), largo demais para o formato.  Aqui: ~118 mm de
% largura, margem interna um pouco maior que a externa (encadernação).
\usepackage[a5paper,inner=16mm,outer=14mm,top=14mm,bottom=16mm,
            includehead,headsep=5mm,footskip=8mm]{geometry}

% Títulos de secção sem justificar (evita linhas estouradas em \Large bold, p. ex.
% «Demonstrabilidade») e um pouco de folga para enunciados de tese em negrito.
\usepackage{sectsty}
\allsectionsfont{\raggedright}
\setlength{\emergencystretch}{1.5em}

% Marca da página impressa do original: \pg{286} imprime [286] no fluxo do texto.
\newcommand{\pg}[1]{\,{\footnotesize\textbf{[#1]}}\,}

% Virada de página do original DENTRO de uma nota (\pgn{422}); discreta, para não quebrar a nota.
\newcommand{\pgn}[1]{{\scriptsize\textbf{[#1]}}}

% Asterisco do original (divisão de assunto, sem título): \asterisco centrado em linha própria,
% também válido no meio do parágrafo ou dentro de nota (\footnote{... \asterisco ...}).
\newcommand{\asterisco}{\par\medskip\centerline{\footnotesize$\ast$}\medskip\noindent}

% Gravura recortada do fac-símile (<figure> do TEI; tools/tei2latex.py): fora de float, no ponto
% exato do original.  \gravura[detalhes]{arquivo.png}{fração da largura do texto}{legenda}:
% legenda principal centrada ("FIGURA VIII. — Côrte…"); os detalhes opcionais (<p> do <figure>:
% "A, hemispherio…; — B, …") vão abaixo, em corpo menor e justificados.  Ambos podem ser vazios.
\usepackage{graphicx}
\newcommand{\gravura}[4][]{\par\medskip\begingroup\centering
  \includegraphics[width=#3\linewidth,height=0.75\textheight,keepaspectratio]{#2}\par
  \if\relax\detokenize{#4}\relax\else\smallskip{\small #4\par}\fi
  \if\relax\detokenize{#1}\relax\else\smallskip{\footnotesize\leftskip=1em\rightskip=1em
    \parfillskip=0pt plus 1fil\noindent #1\par}\fi\endgroup\medskip\noindent}

% Citações latinas de S. Tomás (o OCR erra muito aqui: sempre conferir na imagem).
\newenvironment{latim}{\begin{quote}\itshape}{\end{quote}}

% --- Estrutura do autor (promovido de "Artigo III.tex") -----------------------
% enunciado de tese numerada que corre com o texto (segue-se travessão)
\newcommand{\tese}[2]{\bigskip\noindent\textbf{#1.\quad #2}\;\textemdash\;}
% enunciado de tese numerada em linha própria (seguem-se as subdivisões a), b)...)
\newcommand{\teseB}[2]{\bigskip\noindent\textbf{#1.\quad #2}\par}
% subdivisão a), b), c)... com rubrica em itálico
\newcommand{\sub}[2]{\medskip\noindent#1)\quad\emph{#2},\;\textemdash\;}
% subdivisão a), b), c)... sem rubrica
\newcommand{\subp}[1]{\medskip\noindent#1)\quad}

% --- Notas: no pé (padrão) ou correndo no corpo do texto ---------------------
% Passar  \def\notascorridas{}  no -usepretex faz cada nota correr no corpo, em
% corpo menor, num bloco recuado que interrompe a leitura até acabar.  Ex.:
%   latexmk -lualatex -usepretex='\def\fragmento{vol2/.../01_cap_y.tex}\def\notascorridas{}' tex/_preview.tex
\newif\ifnotascorridas
\ifdefined\notascorridas \notascorridastrue \fi
\ifnotascorridas
  \renewcommand\footnote[1]{%
    \stepcounter{footnote}\textsuperscript{\thefootnote}% marca no ponto do texto
    \par
    {\footnotesize
     \leftskip=2em \rightskip=2em \parindent=1em
     \addvspace{\medskipamount}%
     \noindent\textit{Nota \thefootnote.}\enspace #1\par
     \addvspace{\medskipamount}}%
  }
\fi

% --- Sem notas: \def\semnotas{} no -usepretex suprime todas as notas ---------
% (nem a chamada no texto nem o conteúdo).  Usado por  ./compila.sh --sem-notas .
\ifdefined\semnotas
  \renewcommand\footnote[2][]{}%
  \renewcommand\footnotemark[1][]{}%
  \renewcommand\footnotetext[2][]{}%
\fi

% peças pré-textuais (nó pretextual; p@type em web/obra/esquema.py)
\newcommand{\preTituloObra}[1]{\par\vspace{1em}{\centering\Large\bfseries\scshape #1\par}\vspace{1em}}
\newcommand{\preCentro}[1]{\par{\centering #1\par}}
\newcommand{\preEnderecamento}[1]{\par\vspace{1em}{\centering\small\scshape #1\par}}
\newcommand{\preSaudacao}[1]{\par\medskip\noindent #1\par}
\newcommand{\preDataLocal}[1]{\par\medskip\noindent\emph{#1}\par}
\newcommand{\preFecho}[1]{\par\smallskip{\raggedleft\scshape #1\par}}
}{% Preâmbulo comum ao livro e à pré-visualização de fragmentos (tex/_preview.tex).

% Mancha em A5.  Sem isto vale o padrão da classe book (~105 mm de largura e margem
% externa de ~25 mm em frente e verso), largo demais para o formato.  Aqui: ~118 mm de
% largura, margem interna um pouco maior que a externa (encadernação).
\usepackage[a5paper,inner=16mm,outer=14mm,top=14mm,bottom=16mm,
            includehead,headsep=5mm,footskip=8mm]{geometry}

% Títulos de secção sem justificar (evita linhas estouradas em \Large bold, p. ex.
% «Demonstrabilidade») e um pouco de folga para enunciados de tese em negrito.
\usepackage{sectsty}
\allsectionsfont{\raggedright}
\setlength{\emergencystretch}{1.5em}

% Marca da página impressa do original: \pg{286} imprime [286] no fluxo do texto.
\newcommand{\pg}[1]{\,{\footnotesize\textbf{[#1]}}\,}

% Virada de página do original DENTRO de uma nota (\pgn{422}); discreta, para não quebrar a nota.
\newcommand{\pgn}[1]{{\scriptsize\textbf{[#1]}}}

% Asterisco do original (divisão de assunto, sem título): \asterisco centrado em linha própria,
% também válido no meio do parágrafo ou dentro de nota (\footnote{... \asterisco ...}).
\newcommand{\asterisco}{\par\medskip\centerline{\footnotesize$\ast$}\medskip\noindent}

% Gravura recortada do fac-símile (<figure> do TEI; tools/tei2latex.py): fora de float, no ponto
% exato do original.  \gravura[detalhes]{arquivo.png}{fração da largura do texto}{legenda}:
% legenda principal centrada ("FIGURA VIII. — Côrte…"); os detalhes opcionais (<p> do <figure>:
% "A, hemispherio…; — B, …") vão abaixo, em corpo menor e justificados.  Ambos podem ser vazios.
\usepackage{graphicx}
\newcommand{\gravura}[4][]{\par\medskip\begingroup\centering
  \includegraphics[width=#3\linewidth,height=0.75\textheight,keepaspectratio]{#2}\par
  \if\relax\detokenize{#4}\relax\else\smallskip{\small #4\par}\fi
  \if\relax\detokenize{#1}\relax\else\smallskip{\footnotesize\leftskip=1em\rightskip=1em
    \parfillskip=0pt plus 1fil\noindent #1\par}\fi\endgroup\medskip\noindent}

% Citações latinas de S. Tomás (o OCR erra muito aqui: sempre conferir na imagem).
\newenvironment{latim}{\begin{quote}\itshape}{\end{quote}}

% --- Estrutura do autor (promovido de "Artigo III.tex") -----------------------
% enunciado de tese numerada que corre com o texto (segue-se travessão)
\newcommand{\tese}[2]{\bigskip\noindent\textbf{#1.\quad #2}\;\textemdash\;}
% enunciado de tese numerada em linha própria (seguem-se as subdivisões a), b)...)
\newcommand{\teseB}[2]{\bigskip\noindent\textbf{#1.\quad #2}\par}
% subdivisão a), b), c)... com rubrica em itálico
\newcommand{\sub}[2]{\medskip\noindent#1)\quad\emph{#2},\;\textemdash\;}
% subdivisão a), b), c)... sem rubrica
\newcommand{\subp}[1]{\medskip\noindent#1)\quad}

% --- Notas: no pé (padrão) ou correndo no corpo do texto ---------------------
% Passar  \def\notascorridas{}  no -usepretex faz cada nota correr no corpo, em
% corpo menor, num bloco recuado que interrompe a leitura até acabar.  Ex.:
%   latexmk -lualatex -usepretex='\def\fragmento{vol2/.../01_cap_y.tex}\def\notascorridas{}' tex/_preview.tex
\newif\ifnotascorridas
\ifdefined\notascorridas \notascorridastrue \fi
\ifnotascorridas
  \renewcommand\footnote[1]{%
    \stepcounter{footnote}\textsuperscript{\thefootnote}% marca no ponto do texto
    \par
    {\footnotesize
     \leftskip=2em \rightskip=2em \parindent=1em
     \addvspace{\medskipamount}%
     \noindent\textit{Nota \thefootnote.}\enspace #1\par
     \addvspace{\medskipamount}}%
  }
\fi

% --- Sem notas: \def\semnotas{} no -usepretex suprime todas as notas ---------
% (nem a chamada no texto nem o conteúdo).  Usado por  ./compila.sh --sem-notas .
\ifdefined\semnotas
  \renewcommand\footnote[2][]{}%
  \renewcommand\footnotemark[1][]{}%
  \renewcommand\footnotetext[2][]{}%
\fi

% peças pré-textuais (nó pretextual; p@type em web/obra/esquema.py)
\newcommand{\preTituloObra}[1]{\par\vspace{1em}{\centering\Large\bfseries\scshape #1\par}\vspace{1em}}
\newcommand{\preCentro}[1]{\par{\centering #1\par}}
\newcommand{\preEnderecamento}[1]{\par\vspace{1em}{\centering\small\scshape #1\par}}
\newcommand{\preSaudacao}[1]{\par\medskip\noindent #1\par}
\newcommand{\preDataLocal}[1]{\par\medskip\noindent\emph{#1}\par}
\newcommand{\preFecho}[1]{\par\smallskip{\raggedleft\scshape #1\par}}
}
\def\fimavulso{\end{document}}
\begin{document}
}\fi
%
% e termina com  \csname fimavulso\endcsname .  Dentro do livro ou do _preview, \pg já existe
% e nada disto roda; \fimavulso fica indefinido e o \csname vira \relax.
%
% Motor: o projeto é lualatex (fontspec/polyglossia).  Se o editor chamar pdflatex, cai
% num preâmbulo equivalente com inputenc/babel — serve para conferir, não para o PDF final.
\documentclass[11pt,a5paper]{book}
\ifdefined\directlua
  \usepackage{fontspec}
  \usepackage{polyglossia}
  \setmainlanguage{portuguese}
\else\ifdefined\XeTeXrevision
  \usepackage{fontspec}
  \usepackage{polyglossia}
  \setmainlanguage{portuguese}
\else
  \usepackage[utf8]{inputenc}
  \usepackage[T1]{fontenc}
  \usepackage{lmodern}
  \usepackage[portuguese]{babel}
  \DeclareUnicodeCharacter{2192}{\ensuremath{\rightarrow}}
\fi\fi
% raiz do repo: compilado da raiz ou de volN/<parte>/ (dois níveis abaixo)
\IfFileExists{tex/preambulo.tex}{% Preâmbulo comum ao livro e à pré-visualização de fragmentos (tex/_preview.tex).

% Mancha em A5.  Sem isto vale o padrão da classe book (~105 mm de largura e margem
% externa de ~25 mm em frente e verso), largo demais para o formato.  Aqui: ~118 mm de
% largura, margem interna um pouco maior que a externa (encadernação).
\usepackage[a5paper,inner=16mm,outer=14mm,top=14mm,bottom=16mm,
            includehead,headsep=5mm,footskip=8mm]{geometry}

% Títulos de secção sem justificar (evita linhas estouradas em \Large bold, p. ex.
% «Demonstrabilidade») e um pouco de folga para enunciados de tese em negrito.
\usepackage{sectsty}
\allsectionsfont{\raggedright}
\setlength{\emergencystretch}{1.5em}

% Marca da página impressa do original: \pg{286} imprime [286] no fluxo do texto.
\newcommand{\pg}[1]{\,{\footnotesize\textbf{[#1]}}\,}

% Virada de página do original DENTRO de uma nota (\pgn{422}); discreta, para não quebrar a nota.
\newcommand{\pgn}[1]{{\scriptsize\textbf{[#1]}}}

% Asterisco do original (divisão de assunto, sem título): \asterisco centrado em linha própria,
% também válido no meio do parágrafo ou dentro de nota (\footnote{... \asterisco ...}).
\newcommand{\asterisco}{\par\medskip\centerline{\footnotesize$\ast$}\medskip\noindent}

% Gravura recortada do fac-símile (<figure> do TEI; tools/tei2latex.py): fora de float, no ponto
% exato do original.  \gravura[detalhes]{arquivo.png}{fração da largura do texto}{legenda}:
% legenda principal centrada ("FIGURA VIII. — Côrte…"); os detalhes opcionais (<p> do <figure>:
% "A, hemispherio…; — B, …") vão abaixo, em corpo menor e justificados.  Ambos podem ser vazios.
\usepackage{graphicx}
\newcommand{\gravura}[4][]{\par\medskip\begingroup\centering
  \includegraphics[width=#3\linewidth,height=0.75\textheight,keepaspectratio]{#2}\par
  \if\relax\detokenize{#4}\relax\else\smallskip{\small #4\par}\fi
  \if\relax\detokenize{#1}\relax\else\smallskip{\footnotesize\leftskip=1em\rightskip=1em
    \parfillskip=0pt plus 1fil\noindent #1\par}\fi\endgroup\medskip\noindent}

% Citações latinas de S. Tomás (o OCR erra muito aqui: sempre conferir na imagem).
\newenvironment{latim}{\begin{quote}\itshape}{\end{quote}}

% --- Estrutura do autor (promovido de "Artigo III.tex") -----------------------
% enunciado de tese numerada que corre com o texto (segue-se travessão)
\newcommand{\tese}[2]{\bigskip\noindent\textbf{#1.\quad #2}\;\textemdash\;}
% enunciado de tese numerada em linha própria (seguem-se as subdivisões a), b)...)
\newcommand{\teseB}[2]{\bigskip\noindent\textbf{#1.\quad #2}\par}
% subdivisão a), b), c)... com rubrica em itálico
\newcommand{\sub}[2]{\medskip\noindent#1)\quad\emph{#2},\;\textemdash\;}
% subdivisão a), b), c)... sem rubrica
\newcommand{\subp}[1]{\medskip\noindent#1)\quad}

% --- Notas: no pé (padrão) ou correndo no corpo do texto ---------------------
% Passar  \def\notascorridas{}  no -usepretex faz cada nota correr no corpo, em
% corpo menor, num bloco recuado que interrompe a leitura até acabar.  Ex.:
%   latexmk -lualatex -usepretex='\def\fragmento{vol2/.../01_cap_y.tex}\def\notascorridas{}' tex/_preview.tex
\newif\ifnotascorridas
\ifdefined\notascorridas \notascorridastrue \fi
\ifnotascorridas
  \renewcommand\footnote[1]{%
    \stepcounter{footnote}\textsuperscript{\thefootnote}% marca no ponto do texto
    \par
    {\footnotesize
     \leftskip=2em \rightskip=2em \parindent=1em
     \addvspace{\medskipamount}%
     \noindent\textit{Nota \thefootnote.}\enspace #1\par
     \addvspace{\medskipamount}}%
  }
\fi

% --- Sem notas: \def\semnotas{} no -usepretex suprime todas as notas ---------
% (nem a chamada no texto nem o conteúdo).  Usado por  ./compila.sh --sem-notas .
\ifdefined\semnotas
  \renewcommand\footnote[2][]{}%
  \renewcommand\footnotemark[1][]{}%
  \renewcommand\footnotetext[2][]{}%
\fi

% peças pré-textuais (nó pretextual; p@type em web/obra/esquema.py)
\newcommand{\preTituloObra}[1]{\par\vspace{1em}{\centering\Large\bfseries\scshape #1\par}\vspace{1em}}
\newcommand{\preCentro}[1]{\par{\centering #1\par}}
\newcommand{\preEnderecamento}[1]{\par\vspace{1em}{\centering\small\scshape #1\par}}
\newcommand{\preSaudacao}[1]{\par\medskip\noindent #1\par}
\newcommand{\preDataLocal}[1]{\par\medskip\noindent\emph{#1}\par}
\newcommand{\preFecho}[1]{\par\smallskip{\raggedleft\scshape #1\par}}
}{% Preâmbulo comum ao livro e à pré-visualização de fragmentos (tex/_preview.tex).

% Mancha em A5.  Sem isto vale o padrão da classe book (~105 mm de largura e margem
% externa de ~25 mm em frente e verso), largo demais para o formato.  Aqui: ~118 mm de
% largura, margem interna um pouco maior que a externa (encadernação).
\usepackage[a5paper,inner=16mm,outer=14mm,top=14mm,bottom=16mm,
            includehead,headsep=5mm,footskip=8mm]{geometry}

% Títulos de secção sem justificar (evita linhas estouradas em \Large bold, p. ex.
% «Demonstrabilidade») e um pouco de folga para enunciados de tese em negrito.
\usepackage{sectsty}
\allsectionsfont{\raggedright}
\setlength{\emergencystretch}{1.5em}

% Marca da página impressa do original: \pg{286} imprime [286] no fluxo do texto.
\newcommand{\pg}[1]{\,{\footnotesize\textbf{[#1]}}\,}

% Virada de página do original DENTRO de uma nota (\pgn{422}); discreta, para não quebrar a nota.
\newcommand{\pgn}[1]{{\scriptsize\textbf{[#1]}}}

% Asterisco do original (divisão de assunto, sem título): \asterisco centrado em linha própria,
% também válido no meio do parágrafo ou dentro de nota (\footnote{... \asterisco ...}).
\newcommand{\asterisco}{\par\medskip\centerline{\footnotesize$\ast$}\medskip\noindent}

% Gravura recortada do fac-símile (<figure> do TEI; tools/tei2latex.py): fora de float, no ponto
% exato do original.  \gravura[detalhes]{arquivo.png}{fração da largura do texto}{legenda}:
% legenda principal centrada ("FIGURA VIII. — Côrte…"); os detalhes opcionais (<p> do <figure>:
% "A, hemispherio…; — B, …") vão abaixo, em corpo menor e justificados.  Ambos podem ser vazios.
\usepackage{graphicx}
\newcommand{\gravura}[4][]{\par\medskip\begingroup\centering
  \includegraphics[width=#3\linewidth,height=0.75\textheight,keepaspectratio]{#2}\par
  \if\relax\detokenize{#4}\relax\else\smallskip{\small #4\par}\fi
  \if\relax\detokenize{#1}\relax\else\smallskip{\footnotesize\leftskip=1em\rightskip=1em
    \parfillskip=0pt plus 1fil\noindent #1\par}\fi\endgroup\medskip\noindent}

% Citações latinas de S. Tomás (o OCR erra muito aqui: sempre conferir na imagem).
\newenvironment{latim}{\begin{quote}\itshape}{\end{quote}}

% --- Estrutura do autor (promovido de "Artigo III.tex") -----------------------
% enunciado de tese numerada que corre com o texto (segue-se travessão)
\newcommand{\tese}[2]{\bigskip\noindent\textbf{#1.\quad #2}\;\textemdash\;}
% enunciado de tese numerada em linha própria (seguem-se as subdivisões a), b)...)
\newcommand{\teseB}[2]{\bigskip\noindent\textbf{#1.\quad #2}\par}
% subdivisão a), b), c)... com rubrica em itálico
\newcommand{\sub}[2]{\medskip\noindent#1)\quad\emph{#2},\;\textemdash\;}
% subdivisão a), b), c)... sem rubrica
\newcommand{\subp}[1]{\medskip\noindent#1)\quad}

% --- Notas: no pé (padrão) ou correndo no corpo do texto ---------------------
% Passar  \def\notascorridas{}  no -usepretex faz cada nota correr no corpo, em
% corpo menor, num bloco recuado que interrompe a leitura até acabar.  Ex.:
%   latexmk -lualatex -usepretex='\def\fragmento{vol2/.../01_cap_y.tex}\def\notascorridas{}' tex/_preview.tex
\newif\ifnotascorridas
\ifdefined\notascorridas \notascorridastrue \fi
\ifnotascorridas
  \renewcommand\footnote[1]{%
    \stepcounter{footnote}\textsuperscript{\thefootnote}% marca no ponto do texto
    \par
    {\footnotesize
     \leftskip=2em \rightskip=2em \parindent=1em
     \addvspace{\medskipamount}%
     \noindent\textit{Nota \thefootnote.}\enspace #1\par
     \addvspace{\medskipamount}}%
  }
\fi

% --- Sem notas: \def\semnotas{} no -usepretex suprime todas as notas ---------
% (nem a chamada no texto nem o conteúdo).  Usado por  ./compila.sh --sem-notas .
\ifdefined\semnotas
  \renewcommand\footnote[2][]{}%
  \renewcommand\footnotemark[1][]{}%
  \renewcommand\footnotetext[2][]{}%
\fi

% peças pré-textuais (nó pretextual; p@type em web/obra/esquema.py)
\newcommand{\preTituloObra}[1]{\par\vspace{1em}{\centering\Large\bfseries\scshape #1\par}\vspace{1em}}
\newcommand{\preCentro}[1]{\par{\centering #1\par}}
\newcommand{\preEnderecamento}[1]{\par\vspace{1em}{\centering\small\scshape #1\par}}
\newcommand{\preSaudacao}[1]{\par\medskip\noindent #1\par}
\newcommand{\preDataLocal}[1]{\par\medskip\noindent\emph{#1}\par}
\newcommand{\preFecho}[1]{\par\smallskip{\raggedleft\scshape #1\par}}
}
\def\fimavulso{\end{document}}
\begin{document}
}\fi
% PADRÃO-OURO vol. 1 (5.ª ed., 1897), bloco A — Logica, Secção Primeira, fim do Cap. II e
% começo do Cap. III (Raciocinio e argumentação).  Págs. impressas 62–65
% (vol1/paginas/0062.png … 0065.png).  Transcrito à mão da imagem; o OCR serviu só de rascunho.
% Convenções (as do ouro do vol. 2, vol2/02_parte_theodicea/01_cap_existencia-de-deus.tex):
%   ortographia e erros do original preservados (« seu elementos », « pelo tres juizos »,
%   « cum um só olhar »); \pg{n} só no corpo; cada nota inteira no \footnote do ponto de
%   chamada, inclusive o que transborda para a página seguinte (a nota (2) da p. 65 acaba
%   na p. 66); cabeçalho corrente e número de página ficam de fora.
% Crase: nesta edição é « á » (acento agudo), conferido com zoom nas pp. 62 e 64.

\pg{62}

\teseB{73}{Regras da equivalencia.}
\textsc{Regra} 1.ª --- Duas \emph{contradictorias} tornam-se \emph{equivalentes},
antepondo a particula \emph{negativa} ao sujeito de uma d'ellas; porque, assim, a negação
altera a \emph{quantidade} e a \emph{qualidade} da proposição. As duas
\emph{contradictorias} «~todos os homens são sabios~» e «~alguns homens não são sabios~»
tornam-se \emph{equivalentes}, antepondo a particula \emph{negativa} ao sujeito da
primeira, «~nem todos os homens são sabios~» que equivale a «~alguns homens não são
sabios~».

\textsc{Regra} 2.ª --- Duas \emph{contrarias} tornam-se \emph{equivalentes}, pospondo a
\emph{negação} ao sujeito d'uma d'ellas; assim a particula \emph{negativa} altera só a
\emph{qualidade}. A proposição «~todos os homens são eloquentes~», pospondo «~não~» ao
sujeito «~homens~», torna-se equivalente á sua \emph{contraria}, que é «~todos os homens
não são eloquentes~».

As \emph{sub-contrarias} não admittem \emph{equivalencia}\footnote{As
\emph{sub-contrarias} não admittem \emph{equivalencia}; porque, seja qual for o posto da
particula \emph{negativa}, uma proposição \emph{sub-contraria} não pode reduzir-se á
significação da outra.}.

\chapter*{Capitulo Terceiro\\[0.3em]\large Raciocinio e argumentação}

{\small\noindent\textsc{Summario}: --- Raciocinio, seus elementos e divisão. ---
Argumentação e syllogismo. --- Syllogismo simples e suas regras. --- Figuras e modos do
syllogismo simples. --- Syllogismo composto e sua divisão. --- Inducção e suas especies.
--- Sophisma e suas especies.\par}

\section*{Artigo I\\ Raciocinio, seu elementos e divisão}

\tese{74}{Raciocinio.} \emph{Raciocinio} é a \emph{união} ou \emph{separação de duas ideas
pelo facto de ambas convirem, ou uma convir e outra não, com uma terceira idea, com a qual
foram comparadas}; ex.
\begin{quote}\small
Todo o espirito é immortal;\\
A alma humana é espirito:\\
Logo a alma humana é immortal;
\end{quote}
Comparando entre si duas ideas, muitas vezes não podemos perceber \emph{immediatamente} a
sua conveniencia ou desconveniencia.\pg{63} Em tal caso recorremos a uma terceira idea,
para com ella comparar as duas primeiras. D'esta comparação resulta que --- ou ambas as
ideas convêm com a terceira, --- ou ambas não convêm, --- ou uma convem e outra não. Na
\emph{primeira} hypothese, conclue-se a conveniencia das duas ideas; --- na
\emph{segunda}, não póde concluir-se a sua conveniencia nem a sua desconveniencia; --- na
\emph{terceira}, conclue-se a sua desconveniencia. Assim, não conhecendo immediatamente a
relação entre estas duas ideas «~\emph{alma humana}~» e «~\emph{immortal}~»,
comparamol-as com uma terceira «~\emph{espirito}~», e, pela sua conveniencia com esta,
concluimos que ellas convêm entre si\footnote{A perfeição das substancias intellectuaes
consiste no conhecimento da verdade. Mas ha uma differença neste conhecimento, e a
differença funda-se na propria natureza d'essas substancias. As intelligencias
superiores, que são os Anjos, conhecem sempre a verdade cum um só olhar, sem movimento ou
discurso. As intelligencias inferiores, como é a do homem, não chegam sempre ao perfeito
conhecimento da verdade com um só olhar, mas devem fazer um movimento, um discurso,
passando do conhecido para o desconhecido, e assim descobrem e conhecem explicitamente
muitas verdades, que estavam implicitamente contidas em determinados principios. É, pois,
necessario que a nossa intelligencia estabeleça os principios, e depois os analyse e
resolva, para que se veja o que nelles se achava envolvido. D'esta analyse ou resolução
derivam diversos juizos, dispostos entre si por modo, que um deriva dos outros. Esta
deducção de um juizo de outros juizos é propriamente o \emph{raciocinio} (Cf.
\emph{Sum. Th.}, p. 1, q. 85, a. 5).

Erram, portanto, --- \emph{Wundt}, quando dá o nome de \emph{raciocinio} a toda e qualquer
forma de conhecimento (mesmo á simples \emph{sensação}), --- e \emph{Locke}, quando diz
que o \emph{raciocinio} é uma simples acção \emph{intuitiva}. Pelo \emph{raciocinio}, a
intelligencia não vê uma verdade \emph{na} outra, mas \emph{passa de uma verdade para
outra}. Esta passagem realisa-se, quando a intelligencia não vê immediatamente a relação
entre duas diversas ideas.}.

\tese{75}{Elementos do raciocinio.} Os elementos do \emph{raciocinio} são dois, a saber:
\emph{materia e forma}.

\subp{a} A \emph{materia} é constituida pelas \emph{tres ideas}, que são a materia
\emph{remota}, e pelo \emph{tres juizos}, que são a materia \emph{proxima}.

Das tres ideas uma é \emph{sujeito} da conclusão, outra é \emph{predicado} da mesma
conclusão, a terceira é o termo de \emph{comparação}. --- O sujeito e o predicado
chamam-se \emph{extremos}; o sujeito é o \emph{extremo menor}, e o predicado é o
\emph{extremo maior}; a idea, que é o termo da \emph{comparação}, diz-se
\emph{medio}\footnote{O \emph{sujeito} diz-se extremo \emph{menor}, e o \emph{predicado}
denomina-se extremo \emph{maior}; porque o predicado, referindo-se ao sujeito como o
principio \emph{formal} e \emph{determinante} ao principio \emph{material} e
\emph{determinavel}, tem uma extensão \emph{maior} que a do sujeito. --- O termo de
\emph{comparação} diz-se \emph{medio}; não só porque é o \emph{meio} pelo qual julgamos da
\emph{conveniencia}, ou da \emph{desconveniencia}, do sujeito com o predicado, mas tambem
porque a sua \emph{extensão} é, geralmente, \emph{media} entre a extensão do sujeito e a
do predicado.}.

\pg{64}
As tres ideas comparadas entre si constituem \emph{tres juizos}, --- um formado pela
comparação do \emph{predicado} e da idea \emph{media}, e chama-se \emph{premissa maior};
--- outro, pela comparação do \emph{sujeito} e da idea \emph{media}, e chama-se
\emph{premissa menor}; --- o terceiro, pela comparação do \emph{sujeito} e do
\emph{predicado}, e diz-se \emph{conclusão}. --- As \emph{premissas} formam o
\emph{antecedente}; a \emph{conclusão} é o \emph{consequente}.

\subp{b} A \emph{forma} do \emph{raciocinio} é a disposição das \emph{ideas} e dos
\emph{juizos}, pela qual, admittido o \emph{antecedente}, deve \emph{necessariamente}
admittir-se o \emph{consequente}. --- O \emph{nexo}, que em virtude da \emph{forma} existe
entre o \emph{antecedente} e o \emph{consequente}, chama-se \emph{consequencia}. --- Não
são, pois, uma e a mesma coisa o \emph{consequente} e a \emph{consequencia}: e raciocinios
ha, em que a consequencia é verdadeira e o consequente é falso, --- ou em que a
consequencia é falsa e o consequente é verdadeiro, como nestes dois exemplos:
% No original, em duas colunas: rubrica à esquerda, syllogismo à direita.
\begin{quote}\small
I (\emph{consequencia verdadeira e consequente falso})\\
Todo o homem é racional;\\
Ora a planta é homem;\\
Logo a planta é racional;\\
II (\emph{consequencia falsa e consequente verdadeiro})\\
Todo o animal é sensitivo;\\
Ora o cão é animal;\\
Logo o homem é livre\footnote{O \emph{nexo} entre o \emph{antecedente} e o
\emph{consequente} é tão estreito, que, admittindo-se aquelle, deve admittir-se tambem
este. Com effeito, quando julgamos ser verdadeiro um principio universal,
\emph{implicitamente} julgamos ser tambem verdadeira uma conclusão, contida naquelle
principio. Quando, pois, se deduz a conclusão, deve ella \emph{explicitamente} affirmar ou
negar o que \emph{implicitamente} foi affirmado ou negado no principio; aliás uma e a
mesma coisa seria e deixaria de ser ao mesmo tempo e sob o mesmo aspecto. --- O juizo
ácerca da verdade da \emph{consequencia} pertence á Logica; mas o juizo acerca da verdade
do \emph{consequente} pode pertencer a outras sciencias.}.
\end{quote}

\tese{76}{Divisão do raciocinio.} O raciocinio é \emph{affirmativo}, ou \emph{negativo}.
--- É \emph{affirmativo} se affirma a \emph{conveniencia} do sujeito com o predicado,
porque ambos convêm com a \emph{idea media}. --- É \emph{negativo}, se affirma a
\emph{desconveniencia} do sujeito com o predicado, porque um d'elles não convem com a
\emph{idea media}\footnote{Já demos exemplos do syllogismo \emph{affirmativo}. Um exemplo
do syllogismo \emph{negativo} é o seguinte: \emph{todo o homem é racional; o cão não é
homem; logo o cão não é racional}. A conclusão diz que o predicado \emph{racional} não
compete ao sujeito \emph{cão}, porque este não convem com á \emph{idea media}, que é
\emph{homem}.}.

\pg{65}
\tese{77}{Principios do raciocinio affirmativo e negativo.} O raciocinio
\emph{affirmativo} funda-se no principio de \emph{conveniencia}, que é: \emph{Duas coisas,
cada uma das quaes é egual a uma terceira, são eguaes entre si}. --- O raciocinio
\emph{negativo} funda-se no principio de \emph{desconveniencia}, que é: \emph{Duas coisas,
das quaes uma só é egual a uma terceira, não são eguaes entre si}\footnote{Advirta-se que
duas coisas, eguaes a uma terceira, são eguaes entre si \emph{naquelle sentido e naquella
medida, em que são eguaes a essa terceira coisa}. Esta advertencia é da maxima
importancia. Porque não podemos dizer que são eguaes entre si duas coisas, uma das quaes
é egual á \emph{terceira, considerada sob um aspecto}, e a outra é egual á mesma terceira
coisa, \emph{mas considerada sob outro aspecto}; porque nesse caso o \emph{meio termo} não
é um só, mas multiplo, e o raciocinio é um sophisma. Assim é falso este raciocinio ---
\emph{alguns homens são bons} --- \emph{alguns homens são maus} --- \emph{logo alguns bons
são maus}, --- porque a terceira coisa --- \emph{alguns homens} --- não é considerada
sempre sob um e o mesmo aspecto, mas na premissa \emph{maior} sob um aspecto e na
\emph{menor} sob outro e differente aspecto, pois os homens que são bons são diversos dos
que são maus, e assim a terceira coisa é duplice, e portanto não ha comparação das duas
coisas com uma terceira, e falta a base ao raciocinio. --- Advirta-se tambem que os dois
principios de \emph{conveniencia} (\emph{identidade}) e de \emph{desconveniencia}
(\emph{discrepancia}) se fundam e se resolvem no principio de \emph{contradicção} ---
\emph{uma coisa não pode simultaneamente ser e deixar de ser} (\emph{idem non potest simul
esse et non esse}). Aquelles dois principios são verdadeiros, porque verdadeiro é o
principio de \emph{contradicção}. Por isso este é principio \emph{remoto} do raciocinio, e
os outros dois são principios \emph{proximos}.}.

\section*{Artigo II\\ Argumentação e syllogismo}

\tese{78}{Argumentação.} \emph{Argumentação} é a \emph{expressão verbal do raciocinio}.
--- Por isso, tudo quanto dissermos acerca da \emph{natureza} da \emph{argumentação}, das
suas \emph{propriedades} e \emph{leis}, tudo deve attribuir-se ao \emph{raciocinio}.

\tese{79}{Especies de argumentação.} As especies de \emph{argumentação} são duas: --- uma,
em que de duas proposições \emph{descemos} para uma terceira, como do \emph{todo} para a
\emph{parte}: --- outra, em que de duas proposições \emph{subimos} para uma terceira, como
das \emph{partes} para o \emph{todo}. --- A primeira especie é a argumentação
\emph{deductiva}, chamada \emph{syllogismo}, a segunda é a \emph{inductiva}, e diz-se
\emph{inducção}. --- Agora tratamos do \emph{syllogismo}; mais tarde, da
\emph{inducção}\footnote{O \emph{todo}, de que se falla aqui na argumentação, --- não é o
\emph{integral}, composto de partes \emph{integrantes}, como é, por ex., uma casa composta
das \emph{paredes}, do \emph{tecto}, e do \emph{pavimento}, --- nem é o \emph{potencial},
composto de partes \emph{potenciaes}, como é, por ex., a nossa alma composta das
faculdades de \emph{entender}, \emph{querer}, etc., --- mas é o \emph{todo universal},
composto de partes \emph{subjectivas}, como é, por ex., o \emph{animal}, cujas partes
\emph{subjectivas} são o animal \emph{racional} e o animal \emph{irracional}. Na verdade,
attribuindo-se, na argumentação, a uma \emph{parte} o predicado, que convem ao
\emph{todo}, é necessario que essa \emph{parte} contenha em si o \emph{todo}. Ora isto não
se verifica relativamente ao todo \emph{integral}, porque uma parte \emph{integrante} não
é o todo \emph{integral} (assim o \emph{tecto} não é a \emph{casa}), --- nem ao todo
\emph{potencial}, pois uma parte \emph{potencial} não equivale ao todo \emph{potencial}
(assim o \emph{entendimento} não é toda a \emph{alma}), --- mas verifica-se relativamente
ao todo \emph{universal}, pois que as partes \emph{subjectivas} contêm o todo
\emph{universal} (assim dizemos que o \emph{homem é animal}).}.

\csname fimavulso\endcsname
